%! TeX Program = lualatex
\documentclass[english,aspectratio=169,t,9pt]{beamer}
% useful options:
%   aspectratio=169

\usepackage[utf8]{inputenc}
\usepackage{babel}
\usepackage{fixltx2e}
\usepackage{graphicx}
\usepackage{longtable}
\usepackage{float}
%\usepackage{wrapfig}
\usepackage{tikz}
\usetikzlibrary{shapes.symbols,positioning}
\usepackage{xfrac}
\usepackage{soul}
\usepackage{textcomp}
\usepackage{marvosym}
\usepackage{wasysym}
\usepackage{latexsym}
\usepackage{amssymb}
\usepackage{hyperref}
\usepackage{fontawesome}
\tolerance=1000
\usepackage{listings}
\usepackage[iso]{isodate}

% select theme before using any of its features
\usetheme{Comsys}

\setstretch{1.3}

\title{The Template}
\author[Name of presenter]{First Author, Co-Author}
\authorextra{Advisor(s): Name(s)}
\talktype{Bachelor's/Master's Initial/Final / Seminar Talk}
\eventname{PURPOSE/EVENT}
\date{January 1\textsuperscript{st} 2025}

% Use one of these. Images are automatically scaled to match \pagewidth and cropped at the top
% TikZ code is executed in an [overlay] tikzpicture just above the title
\titletikz{
	\node[anchor=west] at (0,0) {the bottom-left corner};
	\node[anchor=east] at (\pagewidth,0) {the bottom-right corner};
	\node at (0.5\pagewidth,3cm) {more-or-less the middle};
}
\titleimage{figures/default-title-page.jpeg}

\begin{document}

\titleframe

\begin{frame}{\st{Agenda}}
	\vfill
	\singleitem{
		\alert{When giving a presentation at COMSYS, you do not need an agenda slide; \\
your presentation should flow naturally and feature reasonable transitions}
		\begin{itemize}
			\item Only for talks $>$45 min you can consider adding such a slide
			\item If in doubt, talk to you advisor
		\end{itemize}
	}
	\vfill
	\vfill
\end{frame}

\begin{frame}{Rules for COMSYS Slides}
	\begin{itemize}
		\item Use only plain, un-animated slide transitions.
		\item Do not use special animation effects
			\begin{itemize}
				\item Things should pop up ``hard''
				\item Things should not fly around, glitter, sparkle, beep, blink, etc.
				\item (Dis-)Appear effects are fine to support the flow of the presentation.
			\end{itemize}
		\item Don't overload your slides with text.
		\item Use graphs and figures.
		\item Don't leave slides half empty -- balance the content.
			\singleitem{Use paragraph spacing settings if necessary and appropriate. Don't simply hit enter multiple times.}
		\item Try to get a consistent look. Re-use colors, shapes, etc.
	\end{itemize}
\end{frame}

\begin{frame}{Blocks}
	\begin{cols}[1mm]{1,1,1}
		\begin{comsyswhiteblock}[sharp corners,equal height group=mygroup]
			White Block \\
			Just a blue frame around the text
		\end{comsyswhiteblock}
		\begin{comsyswhiteblock}[equal height group=mygroup]
			White Block \\
			Just a blue frame around the text
		\end{comsyswhiteblock}
		\col
		\begin{comsysredblock}[sharpish corners,equal height group=mygroup]
			This is a Alert Block \\
			Alert block
		\end{comsysredblock}
		\begin{comsysredblock}[equal height group=mygroup]
			This is a Alert Block \\
			Alert block
		\end{comsysredblock}
		\col
		\begin{comsysgreenblock}[sharp corners,equal height group=mygroup]
			This is a example block \\
			example block
		\end{comsysgreenblock}
		\begin{comsysgreenblock}[equal height group=mygroup]
			This is a example block \\
			example block
		\end{comsysgreenblock}
	\end{cols}
	\begin{cols}[2mm]{0.25,0.5}
		\begin{actionenv}<2-> % Because of reasons, you can't use action specifiers directly on the blocks :(
		\begin{comsysredblock}[hbox,comsys/bottom cutout]%
			important remarks
		\end{comsysredblock}
		\end{actionenv}
		\begin{comsysyellowblock}[before skip=0pt]
			may use a cutout
		\end{comsysyellowblock}
		\col
		\begin{comsysredblock}[colback=RWTHblau,comsys/bottom cutout=0.7,comsys/bottom cutout dx=5mm,comsys/bottom cutout dy=3mm]
			The blocks accept \texttt{tcolorbox} options.
		\end{comsysredblock}
		\begin{comsysblueblock}[before skip=0pt]
			Cutouts have options too!
		\end{comsysblueblock}
		\begin{tcolorbox}[comsysboxbase,adjusted title=A Title,colframe=RWTHpetrol,coltitle=white]
			The boxes use the \texttt{tcolorbox} package, \\ feel free to create your own
		\end{tcolorbox}
	\end{cols}
\end{frame}

\begin{frame}{Itemize}
	\begin{cols}{1,1,1}
		\begin{itemize}
		\item item
		\item item
		\item []
		\item item
			\begin{itemize}
			\item subitem
				\begin{itemize}
				\item subsubitem
				\end{itemize}
			\end{itemize}
		\end{itemize}
	\col
		\begin{enumerate}
		\item Enumerated item
		\item Enumerated item
		\item[]
		\item Enumerated item
		\item Enumerated item
		\end{enumerate}
	\col
		\setlength{\leftmargini}{\labelsep+\widthof{Item:}+1mm+1.5bp}
		\begin{itemize}
		\item[Item:] \textcolor{black}{\normalfont{custom list}}
		\item[Item:] \textcolor{black}{\normalfont{custom list}}
		\item[Item:] \textcolor{black}{\normalfont{custom list}}
		\end{itemize}
	\end{cols}

	\vskip\bigskipamount

	\begin{cols}{auto}
		\begin{comsysblueblock}
			\begin{itemize}
			\item item
				\begin{itemize}
				\item subitem
					\begin{itemize}
					\item subsubitem
					\end{itemize}
				\end{itemize}
			\end{itemize}
		\end{comsysblueblock}
		\begin{comsysblueblock}
			\begin{enumerate}
			\item item
			\item item
			\end{enumerate}
		\end{comsysblueblock}
	\col
		\begin{comsysgreenblock}
			\begin{itemize}
			\item item
				\begin{itemize}
				\item subitem
					\begin{itemize}
					\item subsubitem
					\end{itemize}
				\end{itemize}
			\end{itemize}
		\end{comsysgreenblock}
		\begin{comsysredblock}
			\begin{itemize}
			\item item
				\begin{itemize}
				\item subitem
					\begin{itemize}
					\item subsubitem
					\end{itemize}
				\end{itemize}
			\end{itemize}
		\end{comsysredblock}
	\col
		\begin{comsyswhiteblock}{}
			\begin{itemize}
			\item item
				\begin{itemize}
				\item subitem
					\begin{itemize}
					\item subsubitem
					\end{itemize}
				\end{itemize}
			\end{itemize}
		\end{comsyswhiteblock}
		\begin{comsysgrayblock}{}
			\begin{itemize}
			\item item
				\begin{itemize}
				\item subitem
					\begin{itemize}
					\item subsubitem
					\end{itemize}
				\end{itemize}
			\end{itemize}
		\end{comsysgrayblock}
	\end{cols}
\end{frame}

\begin{frame}{Further Bars}
	\begin{comsysyellowblock}
		This is a takeaway
		\comsysblocksubtitle{0.8}
		With a subtitle
	\end{comsysyellowblock}
	\begin{comsysyellowblock}[comsys/left icon=?]
		Is this a test question?
	\end{comsysyellowblock}
	\begin{comsysredblock}[comsys/right icon=\faThumbsDown]
		Everything is bad
	\end{comsysredblock}\vspace{-0.75ex}
	\begin{comsysredblock}[comsys/right icon=!]
		We need to warn
	\end{comsysredblock}
	\begin{comsysgreenblock}[comsys/right icon=\faThumbsUp]
		Everything is good
	\end{comsysgreenblock}
	\begin{comsysblueblock}[comsys/left icon=\faLightbulbO]
		Because we know how to do it \\ (we're very good)
		\comsysblocksubtitle{0.5}
		sub
	\end{comsysblueblock}
\end{frame}

\begin{frame}{Headings}
	\heading{Title}
	\begin{itemize}
	\item A
	\item B
	\end{itemize}

	\vspace{\bigskipamount}

	\begin{cols}{0.5,auto}
		\heading{Title on the left}
		\begin{itemize}
		\item A
		\item B
		\end{itemize}
	\col
		\heading{Title on the right}
		\begin{itemize}
		\item A
		\item B
		\end{itemize}
	\end{cols}

	\vspace{\bigskipamount}

	\begin{cols}{1,1,1}
		\heading[fill=RWTHgrun]{Title on the left}
		\begin{itemize}
		\item A
		\item B
		\end{itemize}
	\col
		\heading[fill=RWTHturkis]{Title in the middle}
		\begin{itemize}
		\item A
		\item B
		\end{itemize}
	\col
		\heading[fill=RWTHmagenta]{Title on the right}
		\begin{itemize}
		\item A
		\item B
		\end{itemize}
	\end{cols}
\end{frame}

\begin{frame}{Illustrative Pictures right of Main Text}
	\begin{itemize}
	\item A picture can be inserted which is right aligned
		\begin{itemize}
		\item It is vertically aligned with the position\ilgraphics[height=2em]{logos/rwth}
		\\it was inserted in the source
		\end{itemize}
	\item There are two variants
		\begin{itemize}
		\item A variant based on \textit{{\textbackslash}includegraphics} which takes the same arguments as \textit{{\textbackslash}includegraphics}
		\item An environment for any content
		\begin{illustration}
			\begin{tikzpicture}
				\draw[draw=red] (0,0) circle (1em);
				\draw[draw=blue] (0,0) rectangle(2em,2em);
			\end{tikzpicture}
		\end{illustration}
		\end{itemize}
	\end{itemize}

	\vskip3\bigskipamount

	\begin{cols}{2,1,1,1}
		\begin{comsysblueblock}
			Inside a Block\ilgraphics[height=1em]{logos/rwth}
		\end{comsysblueblock}
		\begin{comsysblueblock}[adjusted title={Block Title\ilgraphics[height=1em]{logos/rwth}}]
			Block Body\ilgraphics[height=1em]{logos/rwth}
		\end{comsysblueblock}
	\col
		Top\ilgraphics[height=1em]{logos/rwth}
	\col
		Above\ilagraphics[height=1em]{logos/rwth}
	\col
		Below\ilbgraphics[height=1em]{logos/rwth}
	\end{cols}
\end{frame}

\section{}

\singlemessageframe{This Frame contains\\only a single Message!}

\begin{frame}{Colors}%
	\setstretch{1.3}
	\begin{cols}{auto}
		\showcolor[HTML]{RWTHblau}\\
		\showcolor[HTML]{RWTHhellblau}\\
		\showcolor[HTML]{RWTHmagenta}\\
		\showcolor[HTML]{RWTHgelb}\\
		\showcolor[HTML]{RWTHpetrol}\\
		\showcolor[HTML]{RWTHturkis}\\
		\showcolor[HTML]{RWTHgrun}\\
		\showcolor[HTML]{RWTHmaigrun}\\
		\showcolor[HTML]{RWTHorange}\\
		\showcolor[HTML]{RWTHrot}\\
		\showcolor[HTML]{RWTHbordeaux}\\
		\showcolor[HTML]{RWTHviolett}\\
		\showcolor[HTML]{RWTHlila}
	\col
		\showcolor[HTML]{RWTHblau!50}\\
		\showcolor[HTML]{RWTHhellblau!50}\\
		\showcolor[HTML]{RWTHmagenta!50}\\
		\showcolor[HTML]{RWTHgelb!50}\\
		\showcolor[HTML]{RWTHpetrol!50}\\
		\showcolor[HTML]{RWTHturkis!50}\\
		\showcolor[HTML]{RWTHgrun!50}\\
		\showcolor[HTML]{RWTHmaigrun!50}\\
		\showcolor[HTML]{RWTHorange!50}\\
		\showcolor[HTML]{RWTHrot!50}\\
		\showcolor[HTML]{RWTHbordeaux!50}\\
		\showcolor[HTML]{RWTHviolett!50}\\
		\showcolor[HTML]{RWTHlila!50}
	\end{cols}
\end{frame}

\begin{frame}{Sizes}%
	\showsizes[cm]%
	\showfonts
\end{frame}

\begin{comsyssplitframe}{Emphasize Something on the Side}{0.45}%
	\begin{itemize}
		\item You can use icons from iconmonstr.com
		\item The \texttt{fontawesome} package also has icons
	\end{itemize}
	\comsysframesplit%
	% we have nee, nww, see, sww, sse, ssw, nne, nnw
	\begin{comsysinvertedblock}[halign=flush left,width=0.8\textwidth,comsys/icon circle radius=1.25em,comsys/nww circled icon={\includegraphics[width=1em]{figures/certificate-white.png}}]
		Using certificates is always a good idea.
	\end{comsysinvertedblock}
	\vfill
	\begin{comsysinvertedblock}[halign=flush right,left skip=0.2\textwidth,comsys/icon circle radius=1.25em,comsys/nee circled icon=\faCloud]
		Ever thought about using them in the cloud? \\ You really should.
	\end{comsysinvertedblock}
	\vfill
	\begin{comsysinvertedblock}[halign=flush left,width=0.8\textwidth,comsys/icon circle radius=1.25em,comsys/nww circled icon={\normalsize\hfuzz=5pt\makebox[0.5em]{\includegraphics[width=1.25em]{figures/computer-white.png}}\raisebox{1ex}{\makebox[0.5em]{\includegraphics[width=1.25em]{figures/certificate-white.png}}}},sharp corners=southeast]
		But you should not use them on too many computers. Even not when...
		% We do dark magic with right/left skip, but it should work intuitively
		\comsysblocksubtitle[right skip=0pt,halign=flush left]{0.65}
		...we want to use a subbox! \\ This is really important!
	\end{comsysinvertedblock}
\end{comsyssplitframe}

\begin{frame}{Conclusion: ``The Template'' comes to an end}
	\singleitem[\hspace{0pt}]{Your talk should feature a dedicated conclusion slide that wraps up the presentation.
	\begin{itemize}
		\item Avoid having a final slide with just your contact information on
		\item Instead add a box ``Thank you for your attention'' to signal the end of your talk (as shown here) \\
			while staying on the conclusion slide with all your awesome findings and takeaways
	\end{itemize}
	}
	\pushdown
	\pause
	\begin{center}
	\begin{comsysinvertedblock}[hbox]
		Thank you for your attention!
	\end{comsysinvertedblock}%
	\end{center}
\end{frame}

\activatevideospacing
\setbeamertemplate{background canvas}[video placeholder]

\begin{frame}{Recording a Video with Picture in Picture}
	\begin{itemize}
	\item You can reserve some space for where your video will appear
		\singleitem{Use {\color{black}\texttt{\textbackslash{}activatevideospacing}} before the title slide}
	\item You can (optionally) insert a placeholder of yourself
	\end{itemize}

	\pushdown

	\begin{videospacing}
	\begin{block}{}
		\medskip
		\centering
		This block does not overlap your video.
		\medskip
	\end{block}
	\end{videospacing}
\end{frame}

\end{document}
